\begin{examplebox}
	\textbf{\textcolor{CExample}{例 1（基础）}}
	\textbf{\textcolor{CExample}{题目：}} 已知有限集合 $A,B$ 满足 $|A|=10,\ |B|=15,\ |A\cap B|=5$，求 $|A\cup B|$。\textbf{\textcolor{CExample}{解题步骤：}}
	\begin{description}[style=nextline, leftmargin=2.8em, labelsep=0.5em, itemsep=0.45em]
		\item[\textbf{第一步（识别公式）：}]
			本题直接应用两集合容斥原理。
			\par\textbf{理由：} 已知两个集合及其交集，求并集，符合公式条件。
		\item[\textbf{第二步（代入数值）：}]
			由容斥原理，$|A\cup B| = |A| + |B| - |A\cap B| = 10 + 15 - 5$。
			\par\textbf{理由：} 将已知数据代入公式，进行简单算术。
		\item[\textbf{第三步（计算出结果）：}]
			$10+15-5 = 20$。
			\par\textbf{理由：} 得出最终结果。
	\end{description}
	\textbf{\textcolor{CExample}{关键结论：}}
	\textbf{答案为 $|A\cup B| = 20$。}

	\medskip
	\textbf{\textcolor{CExample}{例 2（稍变形）}}
	\textbf{\textcolor{CExample}{题目：}} 某班共有50人，其中喜欢篮球的有30人，喜欢足球的有20人，两项都喜欢的有8人。问至少喜欢一种球类的人数。\textbf{\textcolor{CExample}{解题步骤：}}
	\begin{description}[style=nextline, leftmargin=2.8em, labelsep=0.5em, itemsep=0.45em]
		\item[\textbf{第一步（建立集合）：}]
			设集合 $A$ 表示喜欢篮球的人，$B$ 表示喜欢足球的人。则 $|A|=30,\ |B|=20,\ |A\cap B|=8$。
			\par\textbf{理由：} 将实际问题转化为集合语言，明确已知数量。
		\item[\textbf{第二步（套用公式）：}]
			至少喜欢一种的人数即 $|A\cup B|$。由两集合容斥原理：
			\[
				|A\cup B| = |A|+|B| - |A\cap B| = 30+20-8.
			\]
			\par\textbf{理由：} 所求即并集大小，直接应用公式。
		\item[\textbf{第三步（计算结果并作答）：}]
			$30+20-8 = 42$。故至少喜欢一种的有42人。
			\par\textbf{理由：} 得出答案并与总人数比较（$42\le 50$，合理）。
	\end{description}
	\textbf{\textcolor{CExample}{关键结论：}}
	\textbf{至少喜欢一种球类的人数为42人。}

	\medskip
	\textbf{\textcolor{CExample}{例 3（综合应用）}}
	\textbf{\textcolor{CExample}{题目：}} 某年级学生参加课外活动：参加数学竞赛的有32人，物理28人，化学25人；同时参加数学和物理的10人，物理和化学的9人，化学和数学的8人；三科都参加的3人。求至少参加一科竞赛的人数。\textbf{\textcolor{CExample}{解题步骤：}}
	\begin{description}[style=nextline, leftmargin=2.8em, labelsep=0.5em, itemsep=0.45em]
		\item[\textbf{第一步（定义集合与数据）：}]
			设 $A,B,C$ 分别表示参加数学、物理、化学竞赛的集合。依题意： $|A|=32,\ |B|=28,\ |C|=25$； $|A\cap B|=10,\ |B\cap C|=9,\ |C\cap A|=8$； $|A\cap B\cap C|=3$。
			\par\textbf{理由：} 将文字信息转化为集合基数的表示，确认所有需要的量。
		\item[\textbf{第二步（应用三集合公式）：}]
			至少参加一科的人数为 $|A\cup B\cup C|$，由三集合容斥原理：
			\[
				\begin{aligned}
					|A\cup B\cup C| &= |A|+|B|+|C| \\ &\quad - |A\cap B| - |B\cap C| - |C\cap A| \\ &\quad + |A\cap B\cap C|.
			\end{aligned}
			\]
			代入数据：
			\[
				|A\cup B\cup C| = 32 + 28 + 25 - 10 - 9 - 8 + 3.
			\]
			\par\textbf{理由：} 识别出适用三集合公式，并正确代入各项。
		\item[\textbf{第三步（分步计算并作答）：}]
			先求和：$32+28+25 = 85$；再减去两两交集之和：$85 - (10+9+8) = 85 - 27 = 58$；最后加上三交集：$58 + 3 = 61$。故至少参加一科的有61人。
			\par\textbf{理由：} 分步计算避免符号错误，并得到最终结果。
	\end{description}
	\textbf{\textcolor{CExample}{关键结论：}}
	\textbf{至少参加一科竞赛的人数为61人。}
\end{examplebox}
